\section{Amplitudes and conventions}
\label{sec:setup}


\subsection{Theoretical framework}

To study the threshold behaviour of $gg\to AA$ through two loops,
I use the amplitudes calculated in
Ref.~\cite{Bhattacharya:2019bmd} with the effective Lagrangian
of Eq.~\eqref{eq:EL}. The labels Type~I and Type~II follow
that reference and distinguish the two classes of diagrams.
They are unrelated to the Yukawa types of two-Higgs-doublet models.
The Type~II tree diagrams correspond, in the full theory, to
two top quark triangles connected by a gluon, i.e. the
one-particle reducible contribution of
Ref.~\cite{Baglio:2023euv}. Here I consider the contribution
with two $O_G$ insertions and omit the superscript
$\mathrm{II}$ from the amplitudes. Diagrams with one $O_G$
and one $O_J$ insertion first contribute at one loop.
Their coefficients are proportional to the $O_G$--$O_G$
tree coefficients, as discussed in
Sec.~\ref{sec:normalization}.

\begin{figure}[t]
 \centering
 \diagramsfigure
 \caption{(a) Type~II tree diagram in the $t$ channel; the $u$-channel diagram is obtained by exchanging the pseudoscalars. The crossed circles are the effective vertices of $O_G$. (b) One-loop Type~II diagram with loop momentum $k$. Its double parton scattering singularity lies at $k=xp_1$ and $p_3-k=(1-x)p_2$, where the four internal gluons are on shell and collinear with the incoming gluons: each incoming gluon splits into two, which fuse pairwise at the effective vertices. Gaunt and Stirling call this topology the crossed box~\cite{GauntStirling2011}.}
 \label{fig:diagrams}
\end{figure}



\subsection{Kinematics}

The Mandelstam variables are
\begin{align}
 \hat s&=(p_1+p_2)^2,
 &t&=(p_1-p_3)^2,
 &u&=(p_2-p_3)^2,
 \label{eq:mandelstam}
\end{align}
with $\hat s+t+u=2\mA^2$. Following the conventions of Ref.~\cite{Bhattacharya:2019bmd}, I use
\begin{align}
 \hat s&=\mA^2\frac{(1+x)^2}{x},
 &t&=-\mA^2 y,
 &u&=-\mA^2 z,
 \label{eq:xyz-definition}
\end{align}
with $x+x^{-1}=y+z$. With $c=\cos\theta$ in the partonic centre of mass frame,
\begin{align}
 \hat s&=\frac{4\mA^2}{1-\bA^2},
 &t,u&=\mA^2-\frac{\hat s}{2}(1\mp\bA c),
 \label{eq:stu-beta}
\end{align}
and the threshold substitution is
\begin{align}
 x&=\frac{1-\bA}{1+\bA},
 &y,z&=\frac{1+\bA^2\mp2\bA c}{1-\bA^2}.
 \label{eq:threshold-map}
\end{align}
All three variables tend to one as $\bA\to0^+$ at fixed $c$, and the expansion has to follow this correlated limit. The Gram determinant and the transverse momentum are
\begin{equation}
 tu-\mA^4
 =\frac{4\mA^4\bA^2(1-c^2)}{(1-\bA^2)^2},
 \qquad
 p_T^2=\frac{\mA^2\bA^2(1-c^2)}{1-\bA^2}.
 \label{eq:gram-scaling}
\end{equation}
Both vanish at threshold at every angle, and also at fixed $\bA$ when $c\to\pm1$.



\subsection{Tensor decomposition and helicity amplitudes}

With the overall factor $C_G^2$ removed, the amplitude is decomposed as
\begin{equation}
 \mathcal{M}^{\mu\nu}_{ab}
 =\delta_{ab}\left(T_1^{\mu\nu}\mathcal{M}_1+T_2^{\mu\nu}\mathcal{M}_2\right),
 \label{eq:tensor-decomposition}
\end{equation}
with
\begin{align}
 T_1^{\mu\nu}
 &=g^{\mu\nu}-\frac{p_1^\nu p_2^\mu}{p_1\!\cdot p_2},
 \label{eq:tensor-one}\\
 T_2^{\mu\nu}
 &=g^{\mu\nu}+\frac{1}{(p_1\!\cdot p_2)p_T^2}
 \Bigl[
 \mA^2p_2^\mu p_1^\nu
 -2(p_1\!\cdot p_3)p_2^\mu p_3^\nu
 \nonumber\\[-2pt]
 &\hspace{3.2em}
 -2(p_2\!\cdot p_3)p_3^\mu p_1^\nu
 +2(p_1\!\cdot p_2)p_3^\mu p_3^\nu
 \Bigr],
 \label{eq:tensor-two}\\
 p_T^2&=\frac{tu-\mA^4}{\hat s}.
 \label{eq:pt-definition}
\end{align}
For physical gluon polarisations, $T_1$ contributes only to equal helicities ($J_z=0$) and $T_2$ only to opposite helicities ($|J_z|=2$), each with a contraction of unit modulus. Up to phases the helicity amplitudes are~\cite{Glover:1987nx}
\begin{equation}
 \mathcal A_{\pm\pm}=C_G^2\,\mathcal M_1,
 \qquad
 \mathcal A_{\pm\mp}=C_G^2\,\mathcal M_2\,e^{\pm2i\phi},
 \label{eq:helicity}
\end{equation}
where $\phi$ is the azimuth of $\vec p_3$, and $\sum_{\mathrm{hel,col}}|\mathcal A|^2=2(\Nc^2-1)\,C_G^4\,(|\mathcal M_1|^2+|\mathcal M_2|^2)$. The labels 1 and 2 refer to the tensors, not to the diagram classes.

At tree level the two Type~II diagrams give
\begin{align}
 \mathcal M_1^{(0)}&=\frac{i\hat s}{2},
 \nonumber\\
 \mathcal M_2^{(0)}&=\frac{i\mA^2(x-z)(x^2+1)(xz-1)}{4xz(x^2-xz+1)}
 \nonumber\\
 &=2i\mA^2\bA^2(1-c^2)+\Ord(\bA^4),
 \label{eq:born}
\end{align}
in the variables of Eq.~\eqref{eq:xyz-definition}. I obtained Eq.~\eqref{eq:born} from the projections of the original calculation and checked it against a direct evaluation of the two Feynman diagrams with the vertex $iC_G\varepsilon_{\mu\alpha\rho\beta}k_1^\mu k_2^\rho\delta^{ab}$ that follows from Eq.~\eqref{eq:EL} (App.~\ref{app:normalization}). The suppression of $\mathcal M_2^{(0)}$ by $\bA^2$ is the centrifugal barrier of the $|J_z|=2$ state.


\subsection{Finite remainders and normalisation}
\label{sec:normalization}

The coefficients $\mathcal M_i$ of Eq.~\eqref{eq:helicity}, from which $C_G^2$ has been factored out, are expanded as
\begin{equation}
 \mathcal{M}_i
 =\mathcal{M}^{(0)}_i
 +\as\mathcal{M}^{(1)}_i
 +\as^2\mathcal{M}^{(2)}_i+\cdots,
 \label{eq:perturbative-expansion}
\end{equation}
with $\as$ renormalised at $\muR$.
The finite remainders, called the finite part in Ref.~\cite{Bhattacharya:2019bmd}, are defined by the IR subtraction of Refs.~\cite{Bhattacharya:2019bmd,Catani:1998bh},
\begin{align}
 \mathcal M^{(1)}&=2I^{(1)}\mathcal M^{(0)}+\mathcal M^{(1),\Fin},
 \nonumber\\
 \mathcal M^{(2)}&=4I^{(2)}\mathcal M^{(0)}+2I^{(1)}\mathcal M^{(1)}+\mathcal M^{(2),\Fin}.
 \label{eq:catani}
\end{align}
From here on $\mathcal M^{(n)}$ denotes the finite remainder
$\mathcal M^{(n),\Fin}$.

At two loops, I decompose each finite remainder into powers
of the number of light quark flavours,
\begin{equation}
 \Mfull{i}=\Mflav{i}{0}+\nf\,\Mflav{i}{1}+\nf^2\,\Mflav{i}{2}.
 \label{eq:nf-decomposition}
\end{equation}
Each $\Mflav{i}{k}$ is finite, but none is the amplitude of
a separate process: the term linear in $\nf$, for example,
combines the diagrams with a closed quark loop with the terms
of the renormalisation and of the infrared subtraction that
depend on $\nf$. In the rest of the paper, the complete
coefficient means $\Mfull{i}$ at $\nf=\BenchNf$.
Furthermore, the remainders released with Ref.~\cite{Bhattacharya:2019bmd} are normalised differently at different orders: the two-loop remainders carry an overall factor $\FileFactor$ relative to the tree and one-loop amplitudes, and the one-loop remainders are those for $\muR=1$~GeV. In the normalisation of Eq.~\eqref{eq:perturbative-expansion} the remainders are therefore
\begin{align}
 \mathcal M_i^{(1)}&=\mathcal M_{i,\mathrm{publ}}^{(1)}+2\beta_0\ln\frac{\muR^2}{\mathrm{GeV}^2}\,\mathcal M_i^{(0)},
 \nonumber\\
 \Mfull{i}&=\frac{1}{\FileFactor}\,\mathcal M_{i,\mathrm{publ}}^{(2)},
 \label{eq:normalisation}
\end{align}
with $\beta_0=(11\Nc-2\nf)/3$ and masses in GeV in the published expressions. The one-loop term restores the scale dependence required by renormalisation group invariance; with it, the released one and two-loop remainders are consistent with each other at every tested point. Appendix~\ref{app:normalization} gives the evidence. The logarithmic structure studied below, the expansion and its accuracy do not depend on these corrections; the one-loop constant and the finite virtual corrections of Sec.~\ref{sec:rate} do.

The one-loop amplitude with one $O_G$ and one $O_J$ insertion enters at $\Ord(\as^4)$, the order of the two-loop $O_G$--$O_G$ amplitude, since $C_J=\Ord(\as)\,C_G$. It is proportional to the tree amplitude: in the normalisation of Eqs.~(3.6) and (3.7) of Ref.~\cite{Bhattacharya:2019bmd},
\begin{equation}
 \mathcal M^{(1),\Fin}_{i,GJ,\mathrm{publ}}
 = \GJRatio\,\nf\,\mathcal M_i^{(0)},
 \qquad i=1,2,
 \label{eq:gj-proportional}
\end{equation}
at every point of the physical phase space. The proportionality is consistent with the axial anomaly, which relates the divergence of the flavour singlet axial current to the gluonic pseudoscalar operator~\cite{Adler:1969gk,Bell:1969ts,Larin1993}. At $\Ord(\as^4)$ the mixed contribution therefore has the kinematic dependence of the tree amplitude and adds no logarithm of $p_T$.

\subsection{Perturbative order}
\label{sec:counting}

The expansion in Eq.~\eqref{eq:perturbative-expansion} is
written for the coefficients $\mathcal M_i$, with the overall
factor $C_G^2$ removed.
The superscript $(n)$ counts loops of gluons and light quarks in the effective theory; in a two-Higgs-doublet model the top quark loops are contained in the Wilson coefficients.
The perturbative order of a contribution to the cross
section also depends on the powers of $\as$ in the Wilson
coefficients.

Since $C_G\propto\as$, the Type~II tree, one-loop and
two-loop contributions to the full amplitude are of
order $\as^2$, $\as^3$ and $\as^4$, respectively.
In a 2HDM, the leading amplitude is of Type~I and
starts at $\Ord(\as)$, giving a LO cross section
of order $\as^2$. Interference between the leading
Type~I amplitude and the Type~II tree amplitude
contributes at order $\as^3$, that is at NLO.
This is the one-particle reducible contribution to
the virtual corrections in
Refs.~\cite{DawsonDittmaierSpira1998,Baglio:2023euv}.
The Type~II two-loop amplitude first enters through
its interference with the leading Type~I amplitude
at order $\as^5$, or N$^3$LO. The same counting
applies to the corresponding class of amplitudes
for Higgs pairs~\cite{Banerjee:2018lfq}.

The mixed $O_G$--$O_J$ contribution discussed in
Sec.~\ref{sec:normalization} also contributes at
this N$^3$LO. In the heavy top quark theory,
$C_J=\Ord(\as)\,C_G$~\cite{Bhattacharya:2019bmd},
so its one-loop amplitude is of order $\as^4$,
the same as the two-loop $O_G$--$O_G$ amplitude.

When no Type~I amplitude is present, the Type~II
tree amplitude instead gives the LO cross section.
This is the case for an axion-like particle whose
only relevant interaction is its coupling to gluons
through $G\tilde G$. The one-loop and two-loop
terms in Eq.~\eqref{eq:perturbative-expansion}
then provide virtual contributions at NLO and
NNLO, respectively. The present results therefore
supply the $O_G$--$O_G$ part of the two-loop virtual
correction at NNLO for this process. A contact
interaction of dimension six containing two
pseudoscalars and two gluon fields would introduce
an additional Type~I amplitude, as in
Ref.~\cite{Brivio:2026alp}.


\subsection{Finite top quark mass effects}

The threshold expansion and the heavy top quark
approximation involve different limits. Taking
$\bA\to0$ places the pair invariant mass close to
$2\mA$, but imposes no restriction on $\mA/m_t$.
Finite top quark mass effects must therefore be
assessed separately when $A$ is the pseudoscalar of a 2HDM,
at values of $\tan\beta$ for which the top quark loop dominates.
No heavy top quark approximation is involved when
the coupling to $G\tilde G$ is specified directly
as an effective interaction, as in the axion-like
particle case above.

For the Type~II tree amplitude, the dependence on
the top quark mass is known. Each triangle in the
full theory couples a pseudoscalar to one on-shell
gluon and one virtual gluon. The two triangles in
a given channel have the same form factor, so the
contribution of each channel is multiplied, relative
to its heavy top quark limit, by
\begin{equation}
 \left[
 \frac{I_2(\tau_A,\lambda)}{I_2^\infty}
 \right]^2,
 \quad
 \tau_A=\frac{4m_t^2}{\mA^2},
 \quad
 \lambda=\frac{4m_t^2}{t}
 \quad\text{or}\quad
 \frac{4m_t^2}{u}.
 \label{eq:top-tree}
\end{equation}
Here $I_2$ is the loop function that also appears
in $A\to Z\gamma$, given in
Ref.~\cite{Baglio:2023euv}, and $I_2^\infty$
denotes its heavy top quark limit.
At the pair threshold, $t=u=-\mA^2$, so both channels have the same factor: \TopMassRatioA, \TopMassRatioB{} and \TopMassRatioC{} for $\mA=\TopMassMAA$, \TopMassMAB{} and \TopMassMAC~GeV, respectively, with $m_t=\BenchMt$~GeV. These differ from unity by at most about \TopMassPercentMax\%.

For the loop corrections, a partial estimate can
be made for the one-loop DPS logarithm. Its
coefficient involves the amplitudes for
$gg\to A$ with both gluons on shell, as discussed
in Sec.~\ref{sec:dps}. The relevant triangle form
factor is therefore obtained in the limit
$\lambda\to\infty$, rather than at the spacelike
virtuality of the exchanged gluon in the tree
diagrams:
\begin{equation}
 \frac{I_2(\tau_A,\infty)}{I_2^\infty}
 =\tau_A f(\tau_A),
 \qquad
 f(\tau)=\arcsin^2\!\left(\frac{1}{\sqrt{\tau}}\right).
 \label{eq:top-dps}
\end{equation}
If the DPS factorisation applies to the full
theory, with the top quark triangles treated as
hard subgraphs, the coefficient of the one-loop
logarithm is multiplied by $[\tau_A f(\tau_A)]^2$.
This gives \TopMassDPSA, \TopMassDPSB{} and
\TopMassDPSC{} for the same three masses, while
the velocity dependence and the helicity
selection rule would remain unchanged.
Neither this estimate nor the comparison of
tree amplitudes determines the finite top
quark mass corrections to the remaining
one-loop terms or to the complete two-loop
amplitude.
